\section{Executable arithmetic and a complete Lean formalization route}
\label{sec:formalization}

\subsection{What has actually been checked}

The companion \code{lean/GammaIPRA.lean} was compiled with Lean
4.32.0, using only \code{Std}. It contains exact rational complex
arithmetic, finite elementary approximations, Euler products, numerical
budget functions, representations of names, and analytic certificate
interfaces. Fourteen named elementary arithmetic lemmas are proved.
The saved \code{\#print axioms} audit reports no axiom dependencies for
each of these fourteen lemmas. Two additional closed arithmetic examples
are checked by kernel reduction. The file contains no admitted proofs
or axiom declarations.

\emph{This does not constitute a formal proof of gamma analysis.}
The analytic certificate structures are uninstantiated obligations.
In particular, the file does not yet prove that its elementary
truncation lengths imply the corresponding analytic error bounds, and
it does not yet produce \code{PRAAx} derivations of the displayed
estimates. These two substantial remaining tasks are specified below.
The human-readable proofs in the preceding sections are the mathematical
argument; the Lean files are an executable foundation and a precise
formalization outline.

\subsection{The exact rational demonstration and its bracket theorem}

The program \code{code/gamma\_brackets.py} evaluates gamma on positive
rational arguments without floating point. We include the proof of its
particular enclosure method so its certificates have a mathematical
interpretation independent of a numerical gamma library.

For a rational $0<r\le1$ set
\[
 A_n(r)=\frac{n!n^r}{r(r+1)\cdots(r+n)},\qquad
 B_n(r)=\frac{(n-1)!n^r}{r(r+1)\cdots(r+n-1)}.
\]
Then
\begin{equation}
 A_n(r)\le\Gamma(r)\le B_n(r),\qquad
 0\le B_n(r)-A_n(r)\le1/n.
 \label{eq:computational-strip-brackets}
\end{equation}
Here is a short complete proof, also explaining the bracket construction
in the inspected Bohr--Mollerup draft~\cite{bmrepo}. For $0<r\le1$,
the concavity inequality $t^r\le1+r(t-1)$ for $t>0$ follows by
integrating the decreasing derivative $r t^{r-1}$ from $1$ to $t$.
Equivalently for rational $r$ it can be cleared to positive integer
power comparisons. Exact cancellation gives
\[
 \frac{A_{n+1}}{A_n}=
 \frac{n+1}{n+1+r}(1+1/n)^r,
 \qquad
 \frac{B_{n+1}}{B_n}=
 \frac{n}{n+r}(1+1/n)^r.
\]
Concavity at $t=1+1/n$ makes the second ratio at most one.
Concavity at $t=n/(n+1)$, followed by positive reciprocation, gives
\[
 (1+1/n)^r\ge(1-r/(n+1))^{-1}\ge1+r/(n+1),
\]
so the first ratio is at least one. Since
$B_n=(1+r/n)A_n$, the brackets are nested. Also
$A_n\le B_n\le B_1=1/r$, and hence
$B_n-A_n=(r/n)A_n\le1/n$. Their common effective limit is the
Euler limit already identified with gamma, proving
\eqref{eq:computational-strip-brackets}.

For a positive rational $x$, integer division computes
$x=r+m$ with $0<r\le1$ and $m\in\N$. This uses a floor only for
a rational input, not for a named real. Put $S=(r)_m>0$. The recurrence
gives $\Gamma(x)=S\Gamma(r)$. For target width $2^{-p}$ the program
chooses
\[
 n=2^{p+2}\max(1,\ceil S).
\]
Write $r=u/v$ in lowest terms, and write $A_n=c_A n^r$,
$B_n=c_B n^r$ with rational $c_A,c_B>0$. It computes a dyadic root
enclosure of $n^r$ by bounded binary search for the integer $k$ such that
\begin{equation}
 k^v\le n^u2^{vL}<(k+1)^v.\label{eq:root-certificate}
\end{equation}
An initial upper endpoint is a power of two whose exponent is obtained
by integer bit length; that many bisection iterations plus one suffice.
Choose
\[
 L=p+3+\operatorname{bitlength}
                    \bigl(\max(1,\ceil{Sc_B})\bigr).
\]
Then $Sc_B2^{-L}\le2^{-p-3}$, while $S/n\le2^{-p-2}$.
The interval
\[
 \left[Sc_A\frac{k}{2^L},\quad
       Sc_B\frac{k+1}{2^L}\right]
\]
contains $\Gamma(x)$ and has width less than $2^{-p}$.
Every quantity in \eqref{eq:root-certificate} is an integer and is
stored in the evidence file. The verification routine reconstructs the
rational coefficients, checks the integer inequalities and interval
width, and checks that the displayed decimal endpoints round outwards.
It does not test the answer against a floating-point gamma oracle.

The supplied run uses eight rational arguments and $p=8$; for example
it gives the rigorous enclosures
\[
 \Gamma(1/2)\in[1.7718050888,1.7726837525],\qquad
 \Gamma(3/2)\in[0.8859025444,0.8863418763].
\]
These finite certificates check implementation arithmetic. Their
interpretation as gamma enclosures uses the bracket theorem above;
samples are not proofs of universal analytic identities.

\subsection{A fully specified rational-input gamma evaluator}

For the complex construction there is also a direct, deliberately
inefficient, closed primitive recursive evaluator. This makes the
formalization target concrete even though the economical strip program
is preferable for examples.

Supply $R\ge1$, $D\in\N$, $\norm z_1\le R$, and the rational tests
$|z+j|\ge1/(D+1)$ for $0\le j\le2R$. Set
\[
 M=M_R,\quad B=R^2+R,\quad
 L=(D+1)^{2R+1}(2R)!3^{2R^2+2R};
 \qquad\lambda=1/L.
\]
For requested precision $p$, choose
\begin{align*}
 N&=\max\{2R,B,3MBL^2 2^{p+6}\},\\
 A&=R(1+R)^N,\qquad T=R(N+1)+1,\\
 H&=A\bigl((R+1)3^{T+1}+1\bigr),\qquad
 \ell=p+H+2L+12.
\end{align*}
Compute $q=\log N$ using the atanh sum with
$N(\ell+2N+4)$ terms, so its error is at most $2^{-\ell}$.
Compute $E=\exp(-zq)$ by the Taylor sum through degree
$4T+\ell+6$. Form the rational complex number
\[
 U=z\prod_{j=1}^N(1+z/j)\,E,
\]
and return its rational inverse. Every loop length is now given by a
closed primitive recursive expression.

To verify the precision budget, the exact polynomial head has
$1$-norm at most $A$. The propagated logarithm and exponential errors
are at most $H2^{-\ell}\le2^{-p-12}/L^2$.
The lower bound for $Q_N$ is $1/L$, so $U$ is apart from zero by at
least $1/(2L)$. Inversion amplifies the $1$-norm error by at most
$4L^2$. Formula (C15) gives the Euler tail at most $2^{-p-6}$ in
complex modulus, hence at most $2^{-p-5}$ in the $1$-norm. The total
error is less than $2^{-p-4}$. The spare margin gives the regular-name
inequality and accommodates the fixed shifts in the equality convention.
Changing $R,D$ to any larger valid domain certificate yields an equal
name by comparing the two finite Euler approximants with the same
finite-tail bound. No search for a modulus occurs.

\subsection{Definitions needed by the full Lean development}

The host-level interfaces in \code{GammaIPRA.lean} are convenience
types for organizing the proof. They are not asserted to be internal
sorts of \ipra. The following definitions fix the intended completion.

\begin{longtable}{@{}>{\raggedright\arraybackslash}p{0.24\textwidth}>{\raggedright\arraybackslash}p{0.70\textwidth}@{}}
\toprule
Object & Definition and required data\\\midrule\endhead
Rational code & Signed integer numerator and positive natural denominator;
every relation is cleared to integer multiplication and comparison.\\
Complex rational & A pair of rational codes with the operations in
\code{CQ}; squared modulus is $u^2+v^2$, and the $1$-norm is $|u|+|v|$.\\
Fixed PR implementation & A repository \code{PR k} expression and its
defining-equation proofs for the chosen routine. The associated function
symbol is used directly; there is no universal evaluation symbol.\\
Regular name & One fixed PR rational-output function, with the
$2^{-p}+2^{-q}$ Cauchy certificate. \code{NameEq} is equality at every
precision with the stated two-stage slack.\\
Domain data & Integer radius $R$, positive rational margin, and finite
pole-distance tests. For a named input, refine its approximation until
the supplied margin is preserved at half its size.\\
Gamma evaluator & The closed evaluator just specified, or its equivalent
reciprocal-product name. Domain-independent equality is a theorem.\\
Euler constant and zeta values & Harmonic-minus-log with (C2), and finite
reciprocal-power sums with (C17).\\
Analytic disk certificate & Coefficient names, a strictly larger rational
disk, coefficient majorant, finite polynomial degree schedule, and
coefficient-rounding and approximation bounds. This is
\code{DiskAnalyticCertificate}.\\
Derivative certificate & A quadratic remainder bound at rational inputs,
with adjustable rational evaluation precision, as in
\code{DerivativeRemainderCertificate}; higher orders iterate finite
polynomial coefficient operations.\\
Integral certificate & Explicit endpoint-tail functions, a compact mesh
error bound, and the finite Riemann-sum program. The integral name is
their specified diagonal evaluation.\\
Derived functions & $\psi=\Gamma'/\Gamma$ with a denominator margin;
polygamma by (C29); beta by (B1); log-gamma by
\eqref{eq:canonical-loggamma}; incomplete gamma by its two truncated
integrals; Bernoulli polynomials by finite antidifferentiation and zero
mean. All have definitions before any theorem invokes them.\\
\bottomrule
\end{longtable}

The analytic certificate in the Lean file uses a coefficient bound
$|c_n|\le M r^{-n}$ on an outer radius $r$ larger than the working
radius. This is enough to prove geometric convergence on the smaller
disk. It must not be replaced by the assertion that some arbitrary
sequence of polynomials approximates the function: that weaker
assertion would not establish analyticity.

\subsection{The proof dependency order and concrete remaining obligations}

\begin{longtable}{@{}>{\raggedright\arraybackslash}p{0.25\textwidth}>{\raggedright\arraybackslash}p{0.42\textwidth}>{\raggedright\arraybackslash}p{0.26\textwidth}@{}}
\toprule
Planned module & Lemmas to implement & Finite invariant / induction\\\midrule\endhead
\code{RatCode} & Sign, fraction, gcd and division implementations;
norm inequalities and error propagation. & Cleared integer identities;
bounded remainder invariants.\\
\code{Elementary} & Exponential and logarithm tails; positive roots;
sine, cosine, arctangent and $\pi$. & Finite sum or product plus a
geometric/factorial tail term.\\
\code{Names} & Regularity, equivalence, arithmetic operations,
margin-preserving inversion, local continuity extension. & Fix the
output precision first; all stage comparisons are rational.\\
\code{EulerProducts} & (C1)--(C15); instantiate the reciprocal and
gamma names and \code{EulerErrorCertificate}. & Partial sums of
$k^{-2}$ and telescoping product error.\\
\code{AnalyticGamma} & (C16)--(C26); instantiate local analytic and
derivative certificates; simple zeros and poles. & Finite coefficient
convolution, binomial translation, geometric inversion.\\
\code{EulerIntegral} & Mesh lemmas and
\eqref{eq:integral-product-error}; derivative integrals and substitutions.
& Compact telescoping sums, explicit small and large tail budgets.\\
\code{FunctionalRelations} & Sine factorization, (T4), reflection,
Gauss, duplication and beta formulas. & Polynomial division by known
separated roots, finite denominator regrouping, beta recurrence.\\
\code{DifferentialGamma} & (C27)--(C45), differentiated reflection and
multiplication, rational-argument derivative values. & Finite power
coefficients and rational reciprocal tails.\\
\code{RealGamma} & Strong convexity, Bohr--Mollerup, Wendel/Gautschi,
ratio estimates and all critical-point algorithms. & Secant inequalities
and rational brackets, preserving endpoint signs.\\
\code{Stirling} & Finite differences, normalization by duplication,
Robbins, finite Euler--Maclaurin and signed remainder. & Finite
telescoper; for each fixed order, rational Bernoulli polynomial
identities and finite reverse-Chebyshev sums.\\
\code{GammaIntegrals} & Binet, Raabe, Mellin, Gaussian and incomplete
gamma consequences. & Finite rectangles and geometric kernels,
with every tail allocated before exchange.\\
\code{ObjectCertificates} & Reify the preceding finite statements as
\code{Fm PTm}; produce \code{Prf PRAAx} derivations. & Use
\code{PRA.qfInd} only on the encoded rational invariants.\\
\bottomrule
\end{longtable}

The entries describe actual proof obligations. None is silently supplied
by an axiom with a name such as ``gamma correctness.'' The earlier
sections give the required finite identity, error bound, or polynomial
argument for each entry. A proof in Mathlib's completed real or complex
field could be added for exposition, but cannot replace
\code{ObjectCertificates} when the target is provability in iPRA.

\subsection{One explicit object-theory target}

Let $\mathsf A_{R,D}(z,p)$ denote the rational evaluator above.
An internal functional-equation target is the integer-coded formula
\[
 \begin{aligned}
 &\forall R,D,z,p\;\mathrm{Domain}(R,D,z,z+1)\ \longrightarrow\\
 &\quad\left\|
 \mathsf A_{R+1,D}(z+1,p+R+6)
 -z\,\mathsf A_{R+1,D}(z,p+R+6)
 \right\|_1\le2^{-p}.
 \end{aligned}
\]
Here $z$ abbreviates two rational codes; \code{Domain} is the finite
list of squared-distance and denominator comparisons, expanded by a
primitive recursive Boolean conjunction. The displayed inequality is
also a Boolean comparison after clearing positive denominators. The
proof starts from the exact finite relation
\[
 (z+N+1)\Gamma_N(z+1)=Nz\Gamma_N(z),
\]
proved by the finite rising-product identity, and bounds its residual
using $|z+N+1|\ge N/2$ at $N\ge2(R+1)$. Substituting the explicitly
chosen common index and the arithmetic evaluation errors proves the
integer-coded formula by quantifier-free induction and equality rules.
Comparing arbitrary two cutoff indices gives the Cauchy theorem; it
does not require the theory to quantify over a limit object.

Reflection, for example, has exactly the same form with
$\sin(\pi z)\Gamma(z)\Gamma(1-z)-\pi$ on the left and the finite
sine estimate (T4) in place of the recurrence residual. Clearing the
identity before dividing is preferable in the formal code. The integral
statements compare two explicitly defined rational algorithms with a
$2^{-p}$ error, not a bare Lean \code{HasIntegral} assertion.

Once an ordinary arithmetic sequent $h:\Gamma\vdash[\code{PRAAx}]\varphi$
has been built, the inspected repository theorem gives its native
translation directly:
\[
 \code{pra\_iff\_native.mp}\ h
 :\ \code{NPrf}\bigl(\Gamma.\code{map}\ \code{NFm.ofPRA}\bigr)
                    (\code{NFm.ofPRA}\ \varphi).
\]
This is where the existing ConciseIPRA result is reused. The bridge
translates a supplied proof; it does not automatically prove new gamma
estimates from their truth in the host natural numbers.

\subsection{Induction and uniformity audit}

The rational evaluator, the critical-point algorithms, finite coefficient
operations, and all stated cutoff formulas are primitive recursive.
Proofs of finite algebraic recursions use natural-number parameters
internally. Assertions indexed by derivative or asymptotic order may
always be read as \emph{a family of iPRA proofs, one for each fixed
order}. The present outline does not silently strengthen this to
induction over an unrestricted analytic predicate. A future uniform
object theorem in the order parameter requires reifying the finite
polynomial invariant, exactly as specified in the module table.

No proof uses an unbounded search for a zero, excluded middle for an
arbitrary real comparison, countable choice, an arbitrary Cauchy
completeness axiom, or the floor of an arbitrary real. Unknown signs in
root finding use the quantitative third branch of the bisection
algorithm. Every complex division carries an explicit margin; all
global assertions about gamma at its poles use the analytic reciprocal
or a punctured Laurent expression.
